In the first test case of the example $n = 3$. Initially the cache holds keys $1, 2, 3$ with value $0$; key $3$ is the most recently used and key $1$ is the least recently used one.

\begin{itemize}
  \item \texttt{2 2}: key $2$ is present with value $0$ at order $2$, so the answer is \texttt{0 2}. Afterwards the recency order is $(2, 3, 1)$ from most to least recent.
  \item \texttt{2 5}: key $5$ is not in the cache, answer \texttt{-1}.
  \item \texttt{1 4 10}: the cache is full, so the least recently used key $1$ is discarded and key $4$ with value $10$ is inserted; the order becomes $(4, 2, 3)$.
  \item \texttt{2 1}: key $1$ has been invalidated, answer \texttt{-1}.
  \item \texttt{2 4}: key $4$ has value $10$ and order $1$, answer \texttt{10 1}.
  \item \texttt{1 2 7}: key $2$ is already in the cache, so only its value changes to $7$ and it becomes most recently used; order is $(2, 4, 3)$.
  \item \texttt{2 2}: answer \texttt{7 1}.
  \item \texttt{1 9 1}: the least recently used key $3$ is discarded and key $9$ with value $1$ is inserted; order is $(9, 2, 4)$.
  \item \texttt{2 3}: key $3$ was discarded, answer \texttt{-1}.
  \item \texttt{2 9}: answer \texttt{1 1}.
  \item \texttt{1 4 8}: key $4$ is still in the cache: its value becomes $8$ and it becomes most recently used; order is $(4, 9, 2)$.
  \item \texttt{2 2}: key $2$ has value $7$ at order $3$, answer \texttt{7 3}.
\end{itemize}

In the second test case the cache has capacity $1$: inserting key $5$ invalidates key $1$. Note how the first \texttt{2 1} query in the third test case refreshes key $1$, so the subsequent insertion of key $3$ invalidates key $2$ instead.
